\documentclass[10pt,a4paper]{amsart}
\usepackage[margin=19mm]{geometry}
\usepackage{amsmath,amssymb,mathtools}
\usepackage[scr=rsfs,cal=euler]{mathalfa}
\usepackage{xfrac}
\usepackage[unicode,hidelinks,bookmarks=false]{hyperref}
\newcommand{\ie}{i.e.\ }
\newcommand{\cf}{cf.\ }
\newcommand{\resp}{respectively\ }
\newcommand{\Subsection}{\subsection}
\providecommand{\nobreakdash}{\nobreak\hbox{-}}
\setlength{\emergencystretch}{2em}
\multlinegap0pt
\usepackage{amssymb}
\usepackage{xcolor}
\newtheorem{lemma}{Lemma}
\newcommand{\FU}[0]{\mathsf{FU}} % Finite unions (for Hindman's theorem)
\newcommand{\FUT}[0]{\mathsf{FUT}}
\newcommand{\NN}[0]{\mathbb{N}}
\newcommand{\Pfin}{\mathcal{P}_{\mathtt{fin}}}
\title{Hindman's theorem does not code $\emptyset^{(\omega)}$\\ in one application}
\author{Lu Liu, Ludovic Patey}
\date{Source: arXiv:2607.17666v1 (CC BY 4.0)}
\begin{document}
\maketitle

\noindent\emph{Source and licence.} Hindman's theorem does not code $\emptyset^{(\omega)}$\\ in one application. Version arXiv:2607.17666v1 (CC BY 4.0), distributed by arXiv under the Creative Commons Attribution 4.0 International licence, http://creativecommons.org/licenses/by/4.0/, as recorded in the arXiv licence metadata for this exact version and shown on the abstract page https://arxiv.org/abs/2607.17666v1. Full licence notice: \texttt{licenses/arxiv-2607.17666-CC-BY-4.0.txt}.\par\smallskip

\noindent\textit{Editorial scope.} Counted unit: the article's lemma lemma (source0.tex lines 285--288). The article's complete original proof of it is delivered with it (source0.tex lines 289--301, one \texttt{proof} environment of 13 lines). No mathematics is written, repaired or re-derived: the statement and the proof are the article's own lines, in the article's own notation, and no retained line is substituted or re-flowed.  No further source-local context is needed: every cross-reference of every retained line resolves inside the two delivered spans.  Nothing was omitted from any retained span, so the package carries no omission record.  No retained line contains a citation, so the wrapper carries no bibliography and no bibliography entry.  No retained line contains a figure environment, an image-inclusion command or drawing code, so no image is delivered and the package declares no asset dependency.  Editorial formatting only, in addition: an AMS \texttt{amsart} preamble that restates the article's own declarations and macros used by the retained lines, namely the environments \texttt{lemma} on separate counters, together with the standard packages those lines need.  The retained environments print in this document as Lemma 1 in order of appearance, whereas the article prints the same items with the numbering of its own full document.  The complete native source of the article as distributed in the arXiv e-print for this exact version is bundled unchanged as \texttt{sources/205589/source0.tex}.
\begin{lemma}
Let $F \subseteq \Pfin(\NN)$ be a finite set and $X$ an infinite block sequence such that $F$ full-matches~$X$.
There is an infinite block sequence~$H \subseteq F \cup \FU(X)$ such that $H \cap F \neq \emptyset$ and $\FU(H)$ is $f$-monochromatic. 
\end{lemma}

\begin{proof}
Let $\bar f : \Pfin(\NN) \to \ell \times F$ be defined by:
\begin{align}\label{eq:full-match-eq1}
\bar f : b & \mapsto (f(a), a) \notag \\
& \text{such that } f(a) = f(a \cup X[b]) = f(X[b]).
\end{align}
Such an $a \in F$ exists since $F$ full-matches~$X$.
By $\FUT$, there is an infinite block sequence $Y \subseteq \Pfin(\NN)$ such that $\FU(Y)$ is $\bar f$-monochromatic for some color $(\ell^*, a) \in \ell \times F$. In other words, by (\ref{eq:full-match-eq1}), for every~$b \in \FU(Y)$,
\begin{align}\label{eq:full-match-eq2}
f(a) = f(a \cup X[b]) = f(X[b]) = \ell^*
\end{align}
Let $H = \{a\} \cup X[Y]$. For any element~$c \in \FU(H)$, either $c = a$, or $c = X[b]$ for some $b \in \FU(Y)$, or $c = a \cup X[b]$ for some~$b \in \FU(Y)$. In either case, by (\ref{eq:full-match-eq2}), $f(c) = \ell^*$. This completes our proof of the lemma.
\end{proof}

\end{document}
